\section{Introduction}
\IEEEPARstart{최}{최근} 인프라의 급격한 발전과 클라우드 컴퓨팅, 소프트웨어 정의 네트워킹(SDN), 네트워크 기능 가상화(NFV), 그리고 대규모 분산 서비스의 확산은 네트워크 환경의 복잡성을 기하급수적으로 증가시키고 있다. 현대 네트워크는 다양한 프로토콜, 이기종 장비, 동적 트래픽 패턴, 그리고 빈번한 구성 변경을 포함하고 있어 관리 난이도가 지속적으로 증가하고 있다. 기존의 수동 기반 네트워크 관리 방식은 운영자의 경험과 반복적인 수작업에 의존하고 있어 구성 오류 가능성이 높고, 장애 발생 시 신속한 대응이 어려우며, 대규모 네트워크 환경에서 확장성 측면에서도 한계를 가진다\cite{mao2017neural, meng2020interpreting}.

초기 네트워크 자동화 연구는 스크립트 기반 자동화나 정책 기반 관리 기법을 중심으로 발전해 왔으며, 최근에는 머신러닝 기반 접근 방식이 네트워크 이상 탐지, 트래픽 예측, 구성 검증 등 다양한 영역에 적용되고 있다. 특히 자연어 이해와 추론 능력을 갖춘 거대 언어 모델(Large Language Models, LLM)은 복잡한 네트워크 구성 문서를 분석하고, 운영자의 질의를 해석하며, 자동으로 구성 명령을 생성할 수 있는 가능성을 보여주면서 네트워크 관리 분야에서 새로운 패러다임으로 주목받고 있다\cite{liu2025large, wu2024netllm}. LLM은 기존 방법으로는 다루기 어려웠던 비정형 데이터 처리와 코드 생성, 다단계 추론 등에서 높은 성능을 보이며 다양한 실무 태스크로 빠르게 확장되고 있다.~\cite{ref1, ref4, ref27}. 

그러나 기존의 단일 LLM 기반 접근 방식은 네트워크 환경에서 요구되는 높은 신뢰성과 정확성을 보장하는 데 한계를 가진다\cite{mondal2023llm_router, wang2025network_arena}. 특히 네트워크는 외부 도구와의 연동, 동적 환경 인식, 긴 컨텍스트에 걸친 정밀 추론은 단순한 텍스트 생성보다 더 복합적인 조정 과정을 요구하기 때문에 단일 LLM은 복잡한 설정 간 관계나 과정을 완전히 이해하지 못하거나, 근거가 부족한 정보를 생성하는 hallucination 문제를 발생시킬 가능성이 있다\cite{bekri2025bridging, wang2024netconfeval}. 

이러한 한계를 여러 에이전트가 역할을 분담해 태스크를 처리하는 Multi-Agent System(MAS) 접근이 확산되고 있으며~\cite{ref2, arzo2021multiagent, bimo2025intentbased}, LangGraph~\cite{ref3} 등의 프레임워크가 에이전트 간 상태 관리와 도구 호출을 체계화하고 있다.
최근에는 도메인 특화 사전학습~\cite{ref18}, LLM 생성 코드 기반 자동화~\cite{ref19}, 네트워크 태스크 적응 프레임워크~\cite{ref31} 등 다양한 접근이 제안되었고, NetConfEval~\cite{wang2024netconfeval}, NIKA~\cite{wang2025network_arena}, NetPress~\cite{ref7}, 그리고 IETF의 NetConfBench 초안~\cite{ref8}은 평가 방법론을 구체화하고 있다. 이는 네트워크 분야에서도 더 이상 LLM의 가능성 자체를 묻는 단계에 머물지 않고, 어떤 태스크를 어떤 방식으로 평가하고 보완할 것인가가 중요한 문제로 이동하고 있음을 보여준다.

대부분의 기존 벤치마크는 \textbf{지식 검색(Knowledge Retrieval)}에 집중한다. TeleQnA~\cite{ref9}는 3GPP 표준 문서에서 사실을 회상하는지, NetBench~\cite{ref11}는 전문가 큐레이션 지식을 활용하는지를 평가한다. 이는 실제 설정 파일에서 네트워크가 어떻게 동작하는지를 추론하는 \textbf{동작 추론(Behavioral Inference)}과는 별개의 문제이다. \texttt{.cfg} 파일에 ``PE1의 OSPF cost가 10''이라고 명시된 것을 읽는 것(L1)과, 그 설정에 따라 PE1에서 PE2까지의 실제 패킷 경로를 추론하는 것(L4)은 질적으로 구분되는 인지 능력을 요구한다.

이 질문에 체계적으로 답하려면 세 가지 조건이 필요하다. 첫째, 사실 추출부터 경로 추론, 장애 영향 분석까지 인지 난이도를 구분하는 벤치마크가 있어야 한다. 기존 벤치마크는 데이터 플레인 시뮬레이션 능력을 측정하지 않는다. 둘째, 네트워크 데이터에 적합한 평가 지표가 필요하다. BERTScore는 \texttt{10.0.1.1}과 \texttt{10.0.1.10}에 0.9 이상의 유사도를 부여하고, Exact Match는 \texttt{\{r1, r2\}}와 \texttt{\{r2, r1\}}을 오답으로 처리하는 등 구조화된 네트워크 답변에 부적합하다. 셋째, 토폴로지의 크기와 복잡도에 따라 LLM의 해석 능력이 달라지므로, 매번 수작업으로 데이터셋을 구성하는 대신 설정 파일만으로 QA를 자동 생성하는 파이프라인이 요구된다.

본 연구는 하나의 연구 질문을 중심으로 전개된다: \textbf{LLM은 네트워크 설정 파일을 얼마나 정확하게 해석하고 분석할 수 있는가?} 이에 답하기 위해 먼저 벤치마크를 구축하고, 평가 과정에서 드러난 한계를 도구 통합 시스템으로 다루며, 그 결과를 실증적으로 분석한다.

본 연구의 기여는 다음과 같다.
\begin{itemize}
  \item \textbf{C1.} \emph{벤치마크, 자동 생성 파이프라인, 평가 지표.} 실제 설정 파일로부터 5단계 인지 난이도의 QA 쌍을 자동 생성하는 벤치마크 NetConfigQA~2.0을 구축한다. 이중 경로(Dual-Path) 파이프라인은 설정 파일과 정책 파일만 교체하면 토폴로지 규모에 관계없이 데이터셋을 자동 확장하며, 네트워크 데이터의 정답 구조를 반영하는 평가 지표 TA-Acc를 함께 제안한다.
  \item \textbf{C2.} \emph{도구 통합 Multi-Agent 시스템과 기여 분리 실험.} 벤치마크 평가에서 드러난 시뮬레이션 장벽을 다루기 위해 PNETLab, Cisco NSO, Batfish 세 평면을 통합하는 Multi-Agent 시스템 NetAlly를 제안하고, Single LLM $\rightarrow$ Pure MAS $\rightarrow$ NetAlly의 3-way 비교로 구조 효과와 도구 효과를 분리한다.
  \item \textbf{C3.} \emph{실증 분석.} 6개 LLM $\times$ 4개 토폴로지(10\textasciitilde40 노드, 총 9,462 QA)에 대한 평가를 통해 시뮬레이션 장벽을 관찰하고, 성능 변화의 원인을 체계적으로 분석한다.
\end{itemize}

Section~II에서는 관련 연구를 정리한다. Section~III에서는 벤치마크 NetConfigQA~2.0과 TA-Acc를 기술하고, Section~IV에서는 NetAlly 시스템을 설명한다. Section~V에서는 실험 결과를 제시하며, Section~VI에서는 그 의미와 한계를 논의한다. 마지막으로 Section~VII에서 결론을 맺는다.
